\chapter{Label classes at general arity}
\label{ch:chain-general}

We now extend the label-class construction to offspring laws supported on
$\set{0,\dots,J}$. The reduced skeleton may have several arities, and a decorated shape carries
a bouquet of finite bushes at each neck vertex. Conditional shape laws retain the dependence
between the exit arity and its decorations. Transport relabellings remove this dependence from the
label field, while merging adjacent splits identifies the possible presented arities with the
branching semigroup.

\section{The reduced skeleton}

\begin{proposition}[Harris decomposition, general support]
\label{thm:harris-general}
\lean{ChainClasses.transform_expansion}
\leanok
\uses{def:gw}
Let $q\in[0,1)$ be the extinction probability, the root of $s=f(s)$ in $[0,1)$. For a
Galton--Watson tree with offspring law $\theta$ conditioned on infinite diameter, the reduced
skeleton is a Galton--Watson tree with the reduced law $\widetilde\nu$, and the decorations are
independent given the skeleton, each a Galton--Watson tree under the conjugate law.
\end{proposition}

\begin{proof}
\leanok
\uses{thm:harris}
Conditional on a vertex having $j$ children, their descendant trees survive independently with
probability $1-q$. The number of surviving children is therefore binomial with parameters $j$
and $1-q$. Conditioning on at least one surviving child and averaging over $j$ gives the reduced
offspring law $\widetilde\nu$. Given the set of surviving children, the remaining child subtrees
are independent and have the original Galton--Watson law conditioned on extinction, namely the
conjugate law. Iterating this decomposition over the surviving vertices gives the claimed
skeleton and decorations.
\end{proof}

\begin{corollary}[The reduced law has full support]
\label{thm:full-support}
\lean{ChainClasses.reducedLaw_pos}
\leanok
In the bushy regime $\widetilde\nu_k>0$ for every $2\leq k\leq J$.
\end{corollary}

\begin{proof}
\leanok
\uses{thm:harris-general}
Among the terms of the coefficient of $\widetilde\nu_k$, the one with $j=J$ is
$\theta_J\binom{J}{k}(1-q)^kq^{\,J-k}$, and in the bushy regime all four factors are positive:
$\theta_J>0$ since $J$ is the largest offspring number, $q>0$ since $\theta_0>0$, and
$1-q>0$ by supercriticality.
\end{proof}

\begin{definition}[Shapes at general arity]
\label{def:shape-general}
\lean{ChainClasses.GShape}
\leanok
\uses{def:shape}
A \emph{shape} is a pair $\sigma=(m,(\beta_1,\dots,\beta_m))$ with $m\geq1$ and each $\beta_i$ a
finite, possibly empty, list of finite trees with offspring numbers in $\set{0,\dots,J}$. Its
realisation is obtained from a path $v_1\cdots v_m$ by attaching the roots of the trees in
$\beta_i$ by an edge at $v_i$.
\end{definition}

\begin{lemma}[The joint law at a split]
\label{thm:split-joint}
\lean{ChainClasses.splitJoint_not_product}
\leanok
Conditioned on being a split, the arity $k$ of a skeleton vertex and the number $r$ of its
bushes have joint law
\[
 \bbP(k,r)=\frac{\theta_{k+r}\binom{k+r}{k}(1-q)^{k-1}q^{\,r}}{1-\widetilde\theta_1}
 \qquad(k\geq2,\ r\geq0).
\]
For $J\geq3$ this law is not a product of a law in $k$ and a law in $r$.
\end{lemma}

\begin{proof}
\leanok
\uses{thm:harris-general}
The event in question requires a total of $k+r$ children and a choice of exactly $k$ surviving
ones. Before conditioning, its probability is
\[
 \theta_{k+r}\binom{k+r}{k}(1-q)^kq^r.
\]
Passing to the reduced tree contributes the survival normalisation, and conditioning the reduced
vertex to be a split contributes the factor $1-\widetilde\theta_1$. Cancelling the common
survival factor gives the displayed formula. For $J\geq3$, the factor
$\theta_{k+r}\binom{k+r}{k}$ depends jointly on $k$ and $r$, so the conditional law need not
factor into independent arity and bush-count laws.
\end{proof}

The failure of the product form is why the shape law has to be taken conditionally on the
arity.

\begin{definition}[Conditional shape laws]
\label{def:conditional-laws}
\lean{ChainClasses.gMixPMF}
\leanok
\uses{thm:split-joint, def:shape-general}
For $k\in\supp\widetilde\nu$ let $\mu_k$ be the law of the shape at the root of the reduced
skeleton, conditioned on its terminating split having arity $k$.
\end{definition}

\begin{lemma}[Explicit form of the conditional laws]
\label{thm:conditional-explicit}
\lean{ChainClasses.gMixMass_eq_measure}
\leanok
\uses{def:conditional-laws}
Under $\mu_k$ the neck length is geometric, the neck decorations are the conjugate bushes, and,
independently of both, the exit bouquet holds $r$ bushes with probability proportional to
$\theta_{k+r}\binom{k+r}{k}q^{\,r}$, each an independent Galton--Watson tree under the conjugate
law.
\end{lemma}

\begin{proof}
\leanok
\uses{thm:split-joint}
Starting from one split, the successive reduced vertices are unary until the first vertex with at
least two children. Independence of the offspring field makes this waiting time geometric, and
the arity of the terminating split is independent of the waiting time. The decorations along the
neck are conjugate bushes and are independent of both quantities. Conditioning the terminal arity
to equal $k$ therefore changes only its exit bouquet. Summing the joint law of
\cref{thm:split-joint} over the possible values of $r$ gives the stated normalisation and the
conditional bouquet law.
\end{proof}

\begin{lemma}[Conditional independence]
\label{thm:conditional-iid}
\lean{ChainClasses.conditional_iid}
\leanok
\uses{def:conditional-laws}
Conditioned on the arity field $(k(w))_w$ of the reduced skeleton, the shapes $S(w)$ are
independent with $S(w)\sim\mu_{k(w)}$.
\end{lemma}

\begin{proof}
\leanok
\uses{thm:conditional-explicit, thm:harris-general}
Fix a finite collection of reduced-skeleton vertices. The neck leading to each of them is explored
inside a disjoint descendant subtree, so the neck lengths are mutually independent and independent
of the terminal arities. By \cref{thm:harris-general}, the decorations are independent given the
reduced skeleton. Hence, after conditioning on the complete arity field, each shape uses an
independent neck and independent decorations. Its conditional distribution is precisely
$\mu_{k(w)}$ by \cref{thm:conditional-explicit}.
\end{proof}

\begin{lemma}[Uniform mass bounds]
\label{thm:mass-uniform}
\uses{def:conditional-laws}
Every $\mu_k$, $k\in\supp\widetilde\nu$, and the mixture $\mu$ satisfy the two bounds of
\cref{thm:shape-mass} with constants depending on $\theta$ alone.
\end{lemma}

\begin{proof}
\uses{thm:shape-mass, thm:conditional-explicit}
Use the description in \cref{thm:conditional-explicit}. The neck length has a fixed geometric
tail, and every decoration is a total progeny under the bounded, subcritical conjugate law. Relative
to the binary case, the exit bouquet adds at most $J-2$ further such progenies. Their moment
generating functions are finite in one common neighbourhood of the origin. Since only finitely
many arities $k$ occur, the resulting exponential tail is uniform in $k$. A supported shape of
size $n$ is specified by at most a fixed multiple of $n$ local choices from a finite set of
positive probabilities. Taking the smallest of these probabilities gives the uniform pointwise
lower bound.
\end{proof}

\section{Transport and the product form}

\begin{lemma}[Transport relabelling]
\label{thm:relabel}
\lean{ChainClasses.exists_gShapeCoupling_relabel}
\leanok
\uses{def:shape-net}
There is $D_2=D_2(\theta)$ such that for every integer $D\geq D_2$ there are a law $\mu_D$ on
$\cR_D$ and maps $\ell_k\colon\cS\times[0,1)\to\cR_D$ pushing every conditional shape law
forward to the single class law $\mu_D$, with a marked quasi-isometry between a shape and its
class representative and with the label graph potential below the matching threshold.
\end{lemma}

\begin{proof}
\leanok
\uses{thm:conditional-iid, thm:mass-uniform, thm:shape-coupling, thm:dilution}
Choose a reference mixture of the finitely many conditional laws. The uniform estimates of
\cref{thm:mass-uniform} allow the descending cascade used in
\cref{thm:shape-coupling} to couple each $\mu_k$ to this mixture with constants independent of
$k$. Before starting the cascade, fix a shape into the relevant support. On the two alternating
sides, shrink at scale $\floor{\sqrt{D/(216J^2)}}$ so that a cascade pair remains comparable
after projection to the net. Push every coupled law through its representative map and denote the
common image law by $\mu_D$. The construction gives the maps $\ell_k$ and the required marked
comparisons. Finally, \cref{thm:dilution} bounds the number of possible representatives at each
size, and \cref{thm:mass-uniform} supplies their mass bounds. Summing over the sizes places the
potential below the matching threshold for $D\geq D_2$.
\end{proof}

\begin{proposition}[Product form]
\label{thm:product-form}
\lean{ChainClasses.labelMeasure_label_pattern}
\leanok
Fix $D\geq D_2$, and let $(U(w))_w$ be independent uniform variables on $[0,1)$, independent
of the realisation. Label the reduced skeleton by
\[
 x(w)=\ell_{k(w)}(S(w),U(w)).
\]
Then the pairs $(x(w),k(w))$ are i.i.d.\ with the product law
$Q=\mu_D\otimes\widetilde\nu$ on $\cR_D\times\set{2,\dots,J}$. The class $v_0$ of the
one-vertex shape carries $\mu_D(v_0)\geq\tfrac12$.
\end{proposition}

\begin{proof}
\leanok
\uses{thm:relabel, thm:conditional-iid}
Condition on the entire arity field. By \cref{thm:conditional-iid}, the shapes are independent
and $S(w)$ has law $\mu_{k(w)}$. The variables $U(w)$ are also independent, and the defining
property of $\ell_k$ sends each $\mu_k$ to $\mu_D$. Thus, conditionally on the arities, the
labels $x(w)$ are independent with common law $\mu_D$. This conditional law does not depend on
the arity field. Averaging over that field therefore makes $x(w)$ independent of $k(w)$ at each
vertex and makes the pairs i.i.d. with law $\mu_D\otimes\widetilde\nu$.
\end{proof}

\begin{lemma}[Cross-law relabelling]
\label{thm:cross-relabel}
\lean{ChainClasses.exists_gShapeCoupling_cond_mix}
\leanok
There is $D_5=D_5(\theta,\theta')$ such that for every integer $D\geq D_5$ a single class law
$\mu_D$ on $\cR_D$ and maps $\ell_k$ serve both offspring laws at once, again with the marked
quasi-isometry and the potential bound.
\end{lemma}

\begin{proof}
\leanok
\uses{thm:relabel}
Fix the unprimed reference mixture $\mu$ and set
$\mu_D=\mu\circ\rep_D^{-1}$. On the unprimed side, use the maps constructed in
\cref{thm:relabel}. For each conditional law on the primed side, the uniform mass estimates give
a cascade coupling with $\mu$. Use the auxiliary uniform variable to sample the $\mu$-coordinate
of this coupling and then apply $\rep_D$. The resulting primed map also has pushforward
$\mu_D$. The two support-fixing and shrinking maps in the cascade give the marked comparison on
their respective sides, and the same entropy estimate gives the common potential bound.
\end{proof}

\section{Merging and the branching semigroup}

\begin{definition}[Merged shapes]
\label{def:merged-shape}
\lean{ChainClasses.RTree.mergeChild}
\leanok
\uses{def:shape-general}
A \emph{merged shape} is a pair $(\tau,(x_1,\dots,x_k))$ with $\tau$ a finite rooted tree whose
root is the entry and whose exits $x_1,\dots,x_k$ are vertices of $\tau$, repetitions allowed.
The shape at a reduced skeleton vertex of arity $k$ is the merged shape whose tree is its
realisation and whose exit list holds $k$ copies of the exit.
\end{definition}

\begin{lemma}[Merging splits]
\label{thm:merge}
\lean{ChainClasses.RTree.merge_size}
\leanok
\uses{def:merged-shape}
Let $s$ be a split of arity $k_1$ in the reduced skeleton and $t$ one of its children, of arity
$k_2$. Contracting the edge $st$ produces a reduced skeleton in which $s$ has arity
$k_1+k_2-1$, and the contraction is a marked quasi-isometry at a constant depending only on
$\ceil{\log_2 J}$.
\end{lemma}

\begin{proof}
\leanok
The shape at $s$ contributes $k_1$ exits. Absorbing the child split $t$ uses one of those exits
and replaces it by the $k_2$ exits of $t$, leaving $k_1+k_2-1$ exits in total. The two shape
realisations and the edge between them form the tree component of the merged shape. To compare the
old and contracted skeletons, encode the at most $J$ children of any split by distinct binary
words of length $\ceil{\log_2J}$. The contraction changes distances only inside one such bounded
addressing block and moves every exit by the same bounded amount. It is therefore a marked
quasi-isometry with a constant depending only on $\ceil{\log_2J}$.
\end{proof}

\begin{lemma}[Semigroup threshold]
\label{thm:semigroup-threshold}
\lean{ChainClasses.semigroup_threshold}
\leanok
Let $\Lambda$ be the additive submonoid of $\bbN_0$ generated by a nonempty finite set of
positive integers with greatest common divisor $d$. There is $v=v(\Lambda)\geq0$ such that $vd$
and $(v+1)d$ lie in $\Lambda$, and every multiple of $d$ at least $v^2d$ lies in $\Lambda$.
\end{lemma}

\begin{proof}
\leanok
Divide every generator by $d$, so that their greatest common divisor is one. B\'ezout's identity
expresses $1$ as an integer combination of the generators. Move the negative terms to the other
side. The two resulting nonnegative combinations are consecutive integers, say $v$ and $v+1$,
and both belong to the rescaled monoid. Every integer $n\geq v^2$ may be written as
$n=av+b(v+1)$ with $a,b\geq0$, by reducing $n$ modulo $v$. Multiplying back by $d$ gives the
claim.
\end{proof}

\begin{corollary}[The bushy semigroup]
\label{thm:bushy-semigroup}
\lean{ChainClasses.branching_semigroup_eq_top}
\leanok
\uses{def:branching-semigroup}
In the bushy regime $\Lambda_{\widetilde\nu}=\bbN_0$.
\end{corollary}

\begin{proof}
\leanok
\uses{thm:full-support}
Supercriticality forces $J\geq2$, so $\widetilde\nu_2>0$ by \cref{thm:full-support} and the
generator $2-1=1$ lies in $\Lambda_{\widetilde\nu}$. Since $1$ generates $\bbN_0$, the
semigroup is everything.
\end{proof}

\section{The chain regime}

\begin{lemma}[Neck and arity are independent]
\label{thm:chain-independence}
\lean{ChainClasses.chain_joint_factor}
\leanok
Let $\theta_0=0$ and $\theta_1\in(0,1)$, and let $w$ be a vertex of the reduced skeleton, with
$m(w)$ the length of the neck ending at $w$ and $k(w)$ the arity of its split. Then
$\bbP(m(w)=r,\ k(w)=j)=\theta_1^{\,r-1}\theta_j$ for $r\geq1$ and $j\geq2$, a product of its
two marginals.
\end{lemma}

\begin{proof}
\leanok
\uses{thm:harris-general}
Explore the unique line of descent from the split above $w$. Since $\theta_0=0$, every child
subtree is infinite and no reduction changes the offspring numbers. The event
$\set{m(w)=r,k(w)=j}$ consists of exactly $r-1$ successive unary vertices followed by one
vertex of arity $j$. Independence of the offspring counts gives probability
$\theta_1^{r-1}\theta_j$. Summing over $j$ and over $r$ separately shows that this joint law is
the product of the geometric neck law and the arity law conditioned on being at least two.
\end{proof}

\begin{proposition}[Product form in the chain regime]
\label{thm:chain-product-form}
\lean{ChainClasses.chain_product_form}
\leanok
Let $D\geq2$ and label the reduced skeleton by $x(w)=\ell_D(m(w))$. Then the pairs
$(x(w),k(w))$ are i.i.d.\ with the product law $Q=\mu_D\otimes\widetilde\nu$ on
$\bbN_0\times\supp\widetilde\nu$, where $\mu_D(i)=\theta_1^{\,D^i-1}-\theta_1^{\,D^{i+1}-1}$.
\end{proposition}

\begin{proof}
\leanok
\uses{thm:chain-independence, thm:chain-classes}
By \cref{thm:chain-independence}, the unquantised pairs $(m(w),k(w))$ are i.i.d. and their two
coordinates are independent. Applying the deterministic map $\ell_D$ to the first coordinate
preserves both properties. The image of the neck marginal is the law computed in
\cref{thm:chain-classes}, while the second marginal remains $\widetilde\nu$. This gives the
stated product law.
\end{proof}

\section{Cluster and blob presentations}

\begin{definition}[Cluster presentation]
\label{def:cluster}
\lean{ChainClasses.clusterMeasure}
\leanok
\uses{thm:merge}
Fix a revealing depth $R\geq1$ and $\beta\in[0,1]$, and attach an independent uniform variable
$U(w)$ to every vertex of the reduced skeleton. Starting from the root, each cluster root either
stands alone or absorbs its first child, according to $U(w)$ and $\beta$. The cluster's
\emph{presented arity} is the arity of the merged shape, and the next cluster roots are the
splits immediately below.
\end{definition}

\begin{lemma}[The presented pairs]
\label{thm:cluster-law}
\lean{ChainClasses.clusterLaw_tsum}
\leanok
\uses{def:cluster}
Under the cluster presentation the presented neck lengths and presented arities are i.i.d.\
over the presented skeleton, with an explicit arity law supported on the reachable interval.
\end{lemma}

\begin{proof}
\leanok
\uses{thm:chain-product-form, thm:merge}
Explore one cluster at a time. Its decision uses only the root data $(m(w),k(w))$, the auxiliary
variable $U(w)$, and at most the revealed prefix of the first child. Every vertex whose data are
inspected is absorbed into the cluster. The next cluster roots consequently lie at stopping
positions beyond the explored region. By the branching property, their descendant fields are
independent and have the original law. The presented pairs are therefore i.i.d. over the
presented skeleton. If no child is absorbed, the presented arity is $k(w)$. If the first child
is absorbed, it is $k(w)+k(w_1)-1$ by \cref{thm:merge}. Summing these alternatives with their
rule probabilities gives the explicit presented arity law and its reachable support.
\end{proof}

\begin{lemma}[Flattening]
\label{thm:cluster-flat}
\lean{ChainClasses.markedQI_of_collapse}
\leanok
\uses{def:cluster}
A cluster $\set{w,w_1}$ realises a neck of length $m(w)$ whose exits leave at two depths at most
$R$ apart. It admits an $R$-marked quasi-isometry to the flat configuration: a neck of length
$m(w)$ with all $k(w)+k(w_1)-1$ exits at its far end.
\end{lemma}

\begin{proof}
\leanok
\uses{thm:merge}
Keep the outer neck fixed and collapse the inner neck towards the far endpoint of the outer one.
Every exit formerly attached along the inner neck is then regarded as an exit at that endpoint.
The map is onto the flat configuration. It does not increase any distance, while a path loses at
most the length $m(w_1)\leq R$ of the collapsed inner neck. The entry is fixed and every exit
moves by at most $R$, so all metric, density, and mark conditions for an $R$-marked
quasi-isometry hold.
\end{proof}

\begin{definition}[Blob presentation]
\label{def:blob}
\lean{ChainClasses.blobMeasure}
\leanok
\uses{def:cluster}
Fix a revealing depth $R\geq1$, a cap $n\geq1$, and a rule that at each step either stops or
names an unspent exit, as a function of what has been revealed and of an independent uniform
variable. Each cluster root grows a \emph{blob} under this rule, with a presented arity and a
set of next cluster roots.
\end{definition}

\begin{lemma}[The reachable arity laws]
\label{thm:blob-law}
\lean{ChainClasses.mixPow_tsum}
\leanok
\uses{def:blob}
Write $P=\widetilde\nu-1$ for the law of the shift $k-1$ under $\widetilde\nu$. Under the blob
presentation the pairs of presented neck length and presented arity are i.i.d.\ over the
presented skeleton, and the reachable arity laws are the mixtures of convolution powers of $P$
that the rule can realise.
\end{lemma}

\begin{proof}
\leanok
\uses{thm:cluster-law, thm:merge}
Expose a blob together with all data inspected by its rule. Because every inspected vertex is
absorbed, the next cluster roots are stopping positions whose descendant subtrees have not been
examined. The branching property makes those subtrees independent with the original law, proving
the i.i.d. assertion by iteration. If the blob absorbs splits of arities
$k_1,\dots,k_m$, then repeated use of \cref{thm:merge} gives presented arity
$1+\sum_{i=1}^m(k_i-1)$. Conditional on the rule choosing $m$ increments, this sum has law
$P^{*m}$. Averaging over the random stopping rule gives precisely a mixture of convolution
powers of $P$.
\end{proof}

\begin{lemma}[The least shift is common]
\label{thm:least-shift}
\lean{ChainClasses.min_mem_of_generates}
\leanok
\uses{def:branching-semigroup}
Let $\Lambda_{\theta}=\Lambda_{\theta'}=\Lambda$ and let $s_0$ be the least nonzero element of
$\Lambda$. Then $s_0+1$ lies in the support of both reduced laws, and more generally every
minimal generator of $\Lambda$ is a shift of both laws.
\end{lemma}

\begin{proof}
\leanok
Let $A$ be either set of supported shifts. Since $A$ generates $\Lambda$, the element $s_0$
is a finite sum of elements of $A$. It cannot be a sum of two or more nonzero terms, since every
such term lies in $\Lambda$ and is therefore at least $s_0$. Thus $s_0\in A$. The same argument
applies to any generator of $\Lambda$ that is minimal under decomposition as a sum of nonzero
elements. Adding one translates these common shifts back to arities in the supports of both
reduced laws.
\end{proof}

\begin{lemma}[Reachable arities]
\label{thm:reachable-arity}
\lean{ChainClasses.cluster_arity_mem}
\leanok
A cluster whose splits have shifts summing to $s$ has presented arity $s+1$, and $s$ ranges over
$\Lambda_{\theta}$. The arities a chain-regime law can present are therefore
$1+(\Lambda_{\theta}\setminus\set0)$, so two laws with $\Lambda_{\theta}=\Lambda_{\theta'}$ can
present the same set of arities.
\end{lemma}

\begin{proof}
\leanok
\uses{thm:blob-law, thm:least-shift}
Begin with one outgoing slot. Absorbing a split of arity $k_i$ consumes that slot and creates
$k_i$ new ones, increasing the number of available exits by $k_i-1$. A cluster containing splits
of arities $k_1,\dots,k_n$ therefore has
$1+\sum_i(k_i-1)$ exits. The sum belongs to $\Lambda_\theta$ by definition. Conversely, every
$s\in\Lambda_\theta\setminus\set0$ is a finite sum of supported shifts. Absorbing a sequence of
splits with those arities constructs a cluster of presented arity $s+1$.
\end{proof}

\begin{lemma}[Asynchronous common renewal]
\label{thm:common-renewal}
\lean{ChainClasses.finite_common_renewal_completions}
\leanok
Let $P,P'$ be laws with nonempty finite supports $A,B\subset\bbN$ of positive integers. If
$\langle A\rangle=\langle B\rangle=\Lambda$, then two independent renewal walks with increments
$P$ and $P'$ have an asynchronous common renewal time with an exponential tail.
\end{lemma}

\begin{proof}
\leanok
\uses{thm:least-shift}
Run the walks asynchronously: at each step advance a walk whose current total is smaller. The lag
then lies in the finite interval determined by the two largest increments, so it is a finite-state
Markov chain. Equality of the generated monoids implies that from every lag state there is a
finite word of increments that returns the two totals to equality. Since the supports are finite,
we may choose a common length bound $H$ for these correction words and a positive lower bound
$q$ for their probabilities. Conditional on any history, the chance of avoiding equality for the
next $H$ advances is at most $1-q$. Iteration gives
$\bbP(\tau>nH)\leq(1-q)^n$, which is the required exponential tail.
\end{proof}

\begin{lemma}[Matched presentation]
\label{thm:matched-presentation}
\lean{ChainClasses.Matched.matchedShift_window}
\leanok
Let $\theta,\theta'$ be chain-regime laws with $\Lambda_{\theta}=\Lambda_{\theta'}=\Lambda$ and
let $s_0$ be the least nonzero element of $\Lambda$. Then both laws admit blob presentations
whose presented pairs are i.i.d.\ with a common arity law, the presented shifts staying in a
bounded window around the common renewal.
\end{lemma}

\begin{proof}
\leanok
\uses{thm:common-renewal, thm:reachable-arity, thm:blob-law}
Couple the two asynchronous renewal procedures through the rule, without coupling the underlying
offspring fields. Starting at a blob root, the rule advances the lagging side and records the
increments until the common renewal of \cref{thm:common-renewal} is reached. A fixed cap stops the
exploration before it can inspect an unbounded region. An auxiliary uniform variable allocates
the residual mass among the possible stopping values so that the two rules induce the same
presented shift law. The exponential tail of the common renewal makes the missed mass uniformly
small, while all uncapped outcomes lie in a bounded window about the common value. Since each blob
absorbs all inspected vertices, \cref{thm:blob-law} applies independently at every new blob root.
Thus the presented pairs are i.i.d. on both sides and have the same arity marginal.
\end{proof}
